\documentclass[11pt,a4paper]{article}
\usepackage[utf8]{inputenc}
\usepackage[T1]{fontenc}
\usepackage{lmodern}
\usepackage[margin=25mm]{geometry}
\usepackage{amsmath,amssymb,booktabs,array,graphicx,microtype}
\usepackage[numbers,sort&compress]{natbib}
\usepackage[hidelinks]{hyperref}
\usepackage{xcolor}
\newcommand{\pending}[1]{\par\noindent\textcolor{blue!60!black}{\textbf{Results pending.} #1}\par}
\newcommand{\asset}[3]{\begin{figure}[htbp]\centering
  \IfFileExists{#1}{\includegraphics[width=.94\linewidth]{#1}}
  {\fbox{\parbox{.88\linewidth}{\centering\textbf{Figure not yet produced.}\\\small\path{#1}}}}
  \caption{#2}\label{#3}\end{figure}}
\newcommand{\tablepath}[1]{\par\noindent\textit{Machine-readable table to populate this result:} \path{#1}.\par}
\newcommand{\draftstatus}{\begin{center}\small\textbf{Substantive research draft---5 November 2026.}\\
  Numerical results under this protocol have not been produced or validated for this draft.
  Code execution, statistical checks and \LaTeX{} compilation remain to be performed.\end{center}}
\newcommand{\R}{\mathbb{R}}
\newcommand{\ind}{\mathbf{1}}
\author{Th\'eotime Coudray, St\'ephane Goutte}
\date{5 November 2026}

\renewcommand{\draftstatus}{\begin{center}\small
  \textbf{Research draft revised on 11 November 2026.}\\
  Completed 600-scenario outcomes support retrospective, request-specific ceilings.
  No requirement-aware router or criticality reweighting has been trained or evaluated.
  Predictive selection, calibration and online benefit remain prospective.
\end{center}}
\date{11 November 2026}

\title{Routing Energy-Optimization Methods for an Explicit Operational Need\\
Available Time, Economic Cost, Unserved Energy and Selection Overhead}
\begin{document}
\maketitle
\draftstatus
\begin{abstract}
Algorithm selection for energy optimization must specify the service to be
delivered. We separate demonstrated portfolio complementarity from the
prospective ability to identify useful methods before solving. A completed
campaign supplies 16,800 primary trials on 600 scenarios and 1,680 additional
repetitions for direct HiGHS, RH construction with repair, JEPA and an
energy-based sixteen-candidate pipeline. For continuity within thirty
seconds, retrospective coverage is 575/600 for the best fixed method,
593/600 for a family-conditioned fixed choice and 597/600 for an individual
oracle. For five-percent economic quality at 120 seconds, both HiGHS and the
portfolio cover 505/600 cases. Adding a strict unserved-energy cap yields
475 versus 477, with that difference disappearing at a one-MWh cap.
Conditional time improvements remain possible: on the 475 joint-acceptable
HiGHS cases, a perfect selector reduces mean completion time from 36.25 to
34.16 seconds before selection overhead. We propose a router conditioned on
scenario descriptors, deadline, economic tolerance, energy cap, risk tolerance
and cost/time priority. Criticality components and their aggregation are
treated as candidate predictive descriptions, not validated routing thresholds.
Nested grouped learning, calibration, simple family and two-method baselines,
and independent timed evaluation are required. No trained-selector advantage
or reliability guarantee is claimed from the retrospective ceilings.

\end{abstract}

\section{Introduction}
The companion comparison asks which operational services four methods can
deliver. This paper asks whether the method suited to a particular service
can be identified before solving. An operator requesting the lowest possible
cost within thirty seconds without unserved energy poses a different problem
from one requesting one-percent economic quality as quickly as possible.
The latter request does not ensure continuity unless an unserved-energy cap
is also included.

The final campaign makes this distinction concrete. On 100 mathematician
scenarios at thirty seconds, RH returns 96 outputs without ENS versus one
for HiGHS; among 95 valid pairs RH has lower cost on 78. In 120-second trials,
JEPA combines a faster exit than HiGHS, five-percent economic quality and no
ENS on 25 medium-criticality, 43 VRE-with-batteries and 15 network cases.
These are observed useful compromises at the stated allocation, not evidence
that those cases can be recognized before solving. The full-cohort ceilings
below show why a router must target explicit requests and beat simple rules.

Per-instance algorithm selection is established. SATzilla uses
instance-dependent performance models \citep{xu2008}; ASlib provides
algorithm-selection scenarios \citep{bischl2016}; runtime prediction emphasizes
informative features and careful evaluation \citep{hutter2013}. Here the
specific question is selection for a stated energy-service requirement, with
complete process exit, explicit unserved energy and online decision costs.

Changing family-average rankings is not enough. A useful router requires
complementarity under the requested criteria, distinguishable favorable cases
using pre-solve information, and sufficient benefit to pay for selection.
A member need not have an artificial winning niche for the portfolio to be
studied. The strongest fixed choice may remain appropriate for many requests,
including economically easy cases.

\section{Evidence scope and shared protocol}
All 16,800 primary trials are complete on 600 scenarios, 100 per family,
with all four methods and seven independent budgets. Two extra repetitions
on 30 preselected cases bring the total to 18,480. The final acceptance
report records no missing reference and no invalid or late output counted
as successful. Reference optimality is certified on 383/600 cases.

The audited descriptive analysis already exists at
\path{notes/final_results_2026-11-11.json}, with input-file hashes. Its primary
coverage, paired comparisons, threshold sensitivity and retrospective
ceilings are reported here; repetition diagnostics are separate. These
results supersede the historical 126-case narrative. The whole cohort has
now informed profile and narrative choices and is therefore a development
resource, not independent confirmation of any fitted router.

The evidence supports complementarity and its limits. Parameter estimation,
criticality-driven selection, risk calibration and actual online routing
have not been performed. The common protocol fixes the physical and timing
interpretation of the completed outcomes.

\subsection{Physical model and independently checked output}
The study reuses the multilayer flexibility formulation and graph representation
developed in the accompanying project manuscript, while applying the corrected
\texttt{thermal\_transitions\_v1} and \texttt{load\_shifting\_v1} implementations.
The model coordinates thermal generation, nuclear generation, renewable
curtailment, batteries, reservoir and pumped hydro, demand response, line flows
and external trade. Its objective combines operating costs, thermal start-up
costs, flexibility activation, storage cycling, curtailment, unserved energy
and overgeneration penalties. Power-related terms are multiplied by the time
step in hours; start-up costs are charged per event.

Thermal commitment uses exact transition relations, initial conditions and
operating, start-up and shutdown ramp limits. Demand response is lossless load
shifting with a causal energy backlog and terminal repayment; historical prose
describing exponentially disappearing rebound energy does not apply. Storage
tracks efficiencies, retention, initial conditions and terminal bounds. The
network is a capacitated transport model. It does not impose Kirchhoff's laws,
security constraints or individual commitment for every equipment node.
Equipment counts, graph size and the number of free integer decisions are
therefore reported as different quantities.

The shared validator reconstructs physical balances and costs from the complete
dispatch, independently of solver residuals and reported objectives. It checks
finite values, axes, integrality, transitions, ramps, storage, demand-response
repayment, line and trade limits, and nodal balances. The absolute physical and
integrality tolerances are respectively $10^{-3}$ and $10^{-5}$. A feasible
dispatch may include unserved energy because the model permits this penalized
action; the delivered unserved energy is reported separately.

\subsection{Operational success and reference quality}
For scenario $s$, method $m$ and independently allocated budget $b$, let
$F_{smb}$ denote delivery of a physically valid output and normal termination
of the entire process tree within $b$. The measured interval includes scenario
reading, graph construction, online embeddings, candidate generation, repair,
validation, output writing and closure. Static model and library preload is
excluded and reported separately. A forced stop, a late dispatch or an incumbent
owned by a still-running process does not satisfy $F_{smb}$.

A scenario has one common reference cost $C_s^{\rm ref}$, obtained with a
600-second solver limit and 16 CPU threads. We distinguish a certified optimum,
a validated incumbent and a lower bound. A time-limited incumbent is not an
optimum. Training-label computations use their own resource allocation and are
not included in solver timing comparisons. For $|C_s^{\rm ref}|>\epsilon$, define
\begin{equation}
 g_{smb}=\frac{C_{smb}-C_s^{\rm ref}}{|C_s^{\rm ref}|},\qquad
 Y_{smb}(\tau)=F_{smb}\,\ind\{g_{smb}\leq\tau\},\quad \tau=0.05.
 \label{eq:economic}
\end{equation}
Negative gaps are retained: a new method can improve an incumbent reference.
For $|C_s^{\rm ref}|\leq\epsilon=1$ euro, the relative gap is undefined and the
quality condition instead requires $C_{smb}-C_s^{\rm ref}\leq1$ euro. This is an
explicit absolute-tolerance convention, not a relative percentage. Sensitivity
outcomes use $\tau\in\{0.01,0.10\}$, feasibility alone, costs and unserved energy.
If a valid reference is missing, the economic outcome is unknown. An unexecuted
trial is also unknown; neither is encoded as an observed failure.

\subsection{Provenance and statistical units}
The main cohort comprises 600 scenarios, 100 in each of six generator families.
The three persona families and the low, medium and high criticality families
are six distributions, not six ordered severity levels. Four methods are
evaluated at budgets $1,2,5,10,30,60,120$ seconds, giving 16,800 planned primary
trials. Two extra repetitions on 30 preselected cases add 1,680 trials.
Each budget is a new execution; the grid is not a single solver trajectory or
an exact time-to-first-solution observation.

The reserve amendment supersedes identities for affected short budgets while
retaining unchanged 120-second identities, including failures. Analyses reconstruct
the amended schedule rather than enumerate every historical result file.
Source, model, scenario and reference hashes are frozen. Historical files are
read without invoking routines that repair or rewrite checkpoints.
All rows from a scenario, across methods, budgets and repetitions, remain in
one statistical group. A controlled variant shares the group of its base.
Preprocessing is fitted within training groups. The primary coverage analysis
uses the first repetition; repeated trials describe variability separately.
Confidence intervals resample independent scenario groups, not individual rows.


\section{Decision problem: a scenario and an operational request}
\subsection{Inputs, outcomes and acceptance conditions}
Let $\mathcal M=\{\text{HiGHS},\text{RH},\text{JEPA},\text{pipeline}\}$.
Write the request as
\begin{equation}
 r=(b,\tau,\bar E,\alpha,o),\qquad
 \pi:(X_s,r)\longmapsto m\in\mathcal M,
 \label{eq:request}
\end{equation}
where $X_s$ contains pre-solve scenario descriptors and criticality components,
$b$ is the total decision deadline, $\tau$ the relative economic tolerance,
$\bar E$ an unserved-energy cap in MWh, $\alpha$ the tolerated probability of
failure, and $o\in\{\text{cost},\text{time}\}$ the priority among acceptable
methods. The output is one method; the router does not run several portfolio
members and retrospectively choose an output. Each member retains its
existing candidate count, repair procedure and internal cost-based selection.

Let $E_{smd}=\sum_{z,t}P^{\rm uns}_{smdzt}\Delta t$ denote unserved energy for
an output produced with execution allowance $d$. The cost $C_{smd}$ is the
full validated objective, including unserved-energy penalties. These
penalties do not replace the independent service cap. Let $Q_{smd}(\tau)$
be the economic condition in Eq.~\eqref{eq:economic}, with the same convention
for near-zero reference costs. Define
\begin{equation}
 Y_{smd}(\tau,\bar E)
 =F_{smd}\,Q_{smd}(\tau)\,\ind\{E_{smd}\leq\bar E\}.
 \label{eq:joint_requirement}
\end{equation}
Setting $\tau=\infty$ disables the cost threshold and setting $\bar E=\infty$
disables the energy cap; an inactive condition does not require its underlying
measurement. The first case still permits cost minimization as the secondary
objective. A zero-unserved-energy request uses a declared tolerance
$\varepsilon_E=10^{-3}$ MWh, matching the final campaign analysis.
Sensitivity to caps of 1, 10 and 100 MWh is available in the archived analysis.
A confirmatory deployment profile must choose a meaningful cap before its
test outcomes are observed; normalized demand-energy caps would be a distinct
prospective requirement, not an already evaluated result.

An observed failed, late or invalid execution is a failure of the joint
request without an artificial cost or energy value. An unexecuted cell or an
otherwise unassessable quality condition is unknown and its support is
reported. In particular, missing references do not prevent assessment of
a service-only request, but do prevent reference-relative assessment of an
otherwise eligible output.

\subsection{Illustrative request profiles}
Table~\ref{tab:profiles} makes the distinction between constraints and priorities
explicit. Their deterministic acceptance criteria can be evaluated
retrospectively on the completed records; no risk-conditioned selector has
been evaluated. They are exploratory profiles, not preregistered confirmatory
endpoints. The exact profile set, risk tolerances and fallback must be locked
before independent validation. No such evaluation is launched here.

\begin{table}[htbp]
\centering\small
\caption{Service profiles, each crossed with the five independent deadlines
$b\in\{5,10,30,60,120\}$ seconds. Each also specifies $\alpha$.
An infinite cap means the corresponding condition is not requested.}
\label{tab:profiles}
\begin{tabular}{p{.35\linewidth}rrl}
\toprule
Profile & $\tau$ & $\bar E$ & Priority $o$\\
\midrule
Urgent continuity & $\infty$ & $\varepsilon_E$ & Cost\\
Balanced response & 0.05 & $\varepsilon_E$ & Time\\
Economic precision & 0.01 & $\infty$ & Time\\
Precision with continuity & 0.01 & $\varepsilon_E$ & Time\\
\bottomrule
\end{tabular}
\end{table}

The first profile asks for the least expensive response among those predicted
to preserve service within the deadline at the requested reliability.
The third asks for fast economic precision and makes no independent service
promise; the fourth explicitly adds one. Each profile--deadline pair is a
distinct request evaluated from its own executions. The JEPA examples observed
with 120 seconds cannot be relabeled as thirty-second successes.

\subsection{Allowable information and the reference distinction}
The training and evaluation labels may use $C_s^{\rm ref}$, but online inputs
cannot include that reference, its lower bound, its certification status,
a solution, an incumbent or any post-solve statistic. A requested $\tau$ is
a preference input; it is not access to the scenario's unknown reference.
The router predicts the chance of meeting a reference-relative benchmark,
rather than certifying that gap online. If an operator instead provides an
absolute monetary ceiling before solving, that ceiling can define a separate
observable cost condition; it must not be synthesized from a hidden reference.

A requirement for a proof of optimality is separate from cost proximity.
It would require an explicit solver certificate and a different acceptance
label. The current four-action comparison and completed descriptive analysis do not
establish such a certificate-aware router.

Gurobi remains an external direct-solver comparator on its available support,
not a fifth action with selectively observed labels. Scenario-specific
retraining or fine-tuning is outside the routing decision.

\section{Predictions and the selection rule}
\subsection{Joint success, cost and complete execution time}
For each method and execution allowance, the proposed predictor estimates
\begin{align}
 p_m(X_s,d,\tau,\bar E)
 &=\Pr\{Y_{smd}(\tau,\bar E)=1\mid X_s,d,\tau,\bar E\},\\
 c_m(X_s,d,\tau,\bar E)
 &=\mathbb E[C_{smd}\mid Y_{smd}=1,X_s,d,\tau,\bar E],\\
 t_m(X_s,d,\tau,\bar E)
 &=\mathbb E[T_{smd}\mid Y_{smd}=1,X_s,d,\tau,\bar E].
 \label{eq:prediction_targets}
\end{align}
The joint probability must be learned or calibrated as a joint event; separate
marginal probabilities are not multiplied under an untested independence
assumption. Cost and time predictions concern acceptable outputs and are
reported with their support. Failed trials still enter the success model and
the overall evaluation. Their missing cost or natural completion time is not
replaced by zero or by the timeout. In particular, a forced termination does
not provide an exact completion time.

All predictions are conditional on the actual allocation $d$. The grid of
independent budget runs is not a single runtime trajectory, and a short exit
in a longer-budget run is not automatically a time-to-quality observation
for a shorter run. Requests outside the evaluated allowance or requirement
support must be identified as extrapolation.

\subsection{Constrained choice rather than a weighted score}
With estimates denoted by hats, define the supported set
\begin{equation}
 \mathcal A(X_s,r,d)=
 \{m\in\mathcal M:\widehat p_m(X_s,d,\tau,\bar E)\geq1-\alpha\}.
\end{equation}
When this set is nonempty, choose
\begin{equation}
 \pi(X_s,r)=
 \begin{cases}
 \displaystyle\arg\min_{m\in\mathcal A}\widehat c_m,
       & o=\text{cost},\\[3pt]
 \displaystyle\arg\min_{m\in\mathcal A}\widehat t_m,
       & o=\text{time}.
 \end{cases}
 \label{eq:selection}
\end{equation}
Only pre-solve predictions break these comparisons. Remaining ties use the
predicted success probability and then a declared deterministic method order,
never the realized test cost. The priority avoids an arbitrary conversion
between euros, seconds and MWh, while the constraints preserve the requested
economic and service conditions.

If $\mathcal A$ is empty, the request is flagged as unsupported at its stated
risk level. If an output method is still required, a declared fallback selects
the highest predicted joint-success probability and explicitly retains that
flag; it does not silently relax the energy cap, economic threshold or
deadline. Fallback cases and their failures count in the complete evaluation.

The threshold $1-\alpha$ is a model-based decision criterion, not a reliability
guarantee. Calibration, sparse acceptable-output samples and distribution
shift may undermine it. Reliability curves, proper probability scores and
held-out violation rates therefore accompany the selection results.
Even 100 favorable primary outcomes in a family do not by themselves
certify a requested small deployment failure probability. The limited
repetition subset and distribution shift must also be considered.

\section{Complementarity before selector fitting}
\subsection{Requirement-specific ceilings and fixed comparators}
For each complete scenario--request cell, the retrospective coverage ceiling is
\begin{equation}
 O_{sr}=\max_{m\in\mathcal M}Y_{smb}(\tau,\bar E).
\end{equation}
It sees all realized outcomes and ignores selection overhead, so it is not
deployable or an exact bound on unknown expected performance. It also does
not establish a probabilistic risk guarantee from one observed outcome.

Coverage alone can miss useful choice. Where several methods already satisfy
a request, define a secondary oracle outcome
\begin{equation}
 J^*_{sr}=\min_{m:Y_{smb}(\tau,\bar E)=1}J_{smb},
 \qquad
 J_{smb}=
 \begin{cases}
 C_{smb},&o=\text{cost},\\
 T_{smb},&o=\text{time}.
 \end{cases}
\end{equation}
This value is undefined if no method succeeds; those cases remain failures
in coverage. Comparing secondary outcomes only on acceptable cases does not
erase the unsuccessful cases or establish all-scenario superiority. Report
the coverage comparison and the common support of every conditional cost or
time contrast together.

The fixed comparator is selected on development groups separately for each
request, including its deadline and priority. It is fixed across scenarios,
not forced to be the same method for incompatible needs. The choice uses the
same acceptance and cost/time priorities as the router. A family-based rule
must likewise be conditioned on the request. Choosing a comparator on final
test outcomes would contaminate the comparison.

A complete-portfolio diagnostic requires known outcomes for all four methods
under the same request. An observed failure is known, while an unexecuted or
unassessable cell is missing. Report exclusions and their family/budget
distribution; a ceiling on complete cases does not resolve selection bias
from missing observations.

\subsection{Observed ceilings in the completed cohort}
The primary grid contains no missing scenario--method--budget cells.
Table~\ref{tab:final_oracles} therefore compares the same 600 scenarios in every
row. For this descriptive diagnostic, both fixed baselines are chosen
retrospectively on the same cohort. Future validation must choose them using
development data only; these ceilings are not held-out selector scores.

\begin{table}[htbp]\centering\small
\caption{Retrospective request satisfaction out of 600 scenarios. No selection
overhead or probabilistic reliability guarantee is included.}
\label{tab:final_oracles}
\begin{tabular}{p{.39\linewidth}rrr}\toprule
Request & Best fixed & Fixed by family & Individual oracle\\\midrule
No ENS, 30 seconds & 575 & 593 & 597\\
Five-percent quality, 120 seconds & 505 & 505 & 505\\
Five-percent quality and no ENS, 10 seconds & 168 & 168 & 182\\
Five-percent quality and no ENS, 120 seconds & 475 & 475 & 477\\\bottomrule
\end{tabular}
\end{table}

At thirty seconds RH is the best fixed continuity choice, but a family rule
already recovers 18 of the 22 additional cases accessible to the full oracle.
At 120 seconds, the continuity counts are 577, 596 and 598 respectively.
A sophisticated selector must therefore justify its incremental benefit
over family-conditioned decisions, including cost and time at preserved service.

For the joint requirement at ten seconds, VRE-with-batteries cases account
for 13 of the 14 additions over the best fixed choice: HiGHS covers 46/100,
the portfolio 59/100, and six cases are covered only by JEPA. This is a
specific development target for identifying within-family complementarity.
By contrast, the two joint additions at 120 seconds are due to HiGHS ENS
of approximately 0.121 and 0.682 MWh. At a one-MWh cap, HiGHS and the oracle
both cover 477/600. This sensitivity precludes presenting the strict-cap
gain as a broad, threshold-independent routing advantage.

Among the 475 scenarios where HiGHS already satisfies the strict joint
requirement at 120 seconds, the fastest acceptable four-method oracle reduces
mean exit time from 36.25 to 34.16 seconds (about 5.8\%). A HiGHS/RH-only
oracle reaches 36.08 seconds on that same support. There are 68 cases with
at least one second and ten percent saved; their median saving is 5.9 seconds.
On the comparable family supports, mean time changes from 34.60 to 31.16
seconds for medium criticality ($n=100$), 10.99 to 9.47 for VRE/batteries
($n=100$), and 63.48 to 58.81 for network cases ($n=95$). These gains are
conditional, retrospective and measured before feature or selection overhead.

Repetitions also limit single-run labels. At 120 seconds, the strict joint
criterion changes across three repetitions for six of the 30 repeated cases
with HiGHS, three with the pipeline, one with JEPA and none with RH. HiGHS
validity itself is stable on all 30. A stable low-quality method is not
necessarily useful, and a single favorable label is not a calibrated
probability of satisfying a service request.

\subsection{Continuation must match the operational claim}
The first service-v2 diagnostic pooled the seven historical deadlines within
each profile and did not authorize fitting. Following inspection of that
development result, the protocol was amended to evaluate the five deadlines
$5,10,30,60,120$ seconds separately. This is a disclosed development amendment,
not a preregistered confirmatory decision or evidence for a fitted selector.

Each of the twenty profile--deadline pairs receives its own continuation
decision. The coverage route requires an oracle gain of at least two
percentage points and a positive lower endpoint of a grouped 95\% interval.
The secondary route requires a conditional relative-cost gain of at least
one percent or a time saving of at least 0.5 seconds, a positive lower
endpoint on common acceptable pairs, oracle satisfaction at least
$1-\alpha$, and a coverage-gain lower endpoint no smaller than minus one
percentage point. Both routes require at least 90\% known portfolio cells
and twenty scenario groups; the secondary route also requires twenty groups
with acceptable pairs. These numerical thresholds are unchanged.
No average across deadlines or positive result at another deadline can
authorize a request. The intervals are pointwise, not simultaneous
guarantees across the twenty comparisons.

Predictors may share training observations with deadline as an input, while
model choice, coverage, conditional cost/time contrasts and calibration
reports remain separate for each profile--deadline pair. The amended
configuration must be frozen before independent confirmation. No amended
gate result or fitted selector is claimed by this implementation update.

If neither form of useful complementarity is supported, the appropriate result
is a limited or negative routing opportunity for that request. In particular,
a saturated success rate from a fixed method leaves no coverage gain but may
still leave cost or time variation; conversely, a lower conditional cost
cannot justify a selector that loses the demanded service reliability.

\section{From criticality sensitivity to router parameters: a proposed bridge}
\subsection{Keep score parameters separate from operational preferences}
The companion index study distinguishes raw scenario descriptors from the
parameters used to aggregate them. Let $z_j$ be a fixed-bound normalized
component and write
\begin{equation}
 I_{\theta}(s)=\lambda\sum_{j\in\mathcal A}w^A_jz_{sj}
 +(1-\lambda)\sum_{j\in\mathcal B}w^B_jz_{sj},
 \qquad \sum w^A_j=\sum w^B_j=1,\quad w_j\geq0.
 \label{eq:router_index}
\end{equation}
The original common description uses $\lambda=0.5$ and equal within-block
weights. The code calls this mixing parameter \texttt{alpha}; here
$\lambda$ distinguishes it from the request's failure tolerance $\alpha$.
Bounds, weights and clipping are score parameters; $(b,\tau,\bar E,\alpha,o)$
are operational requirements. Changing the former does not change a scenario,
its stored solver outcomes, or the latter's acceptance criteria.

Raw arithmetic contributions explain why a score changes, not why a method
succeeds. A constant component can contribute to the level of $I$ without
distinguishing cases, and a clipped component can hide increasingly extreme
inputs. The implementation fixes minimum-generation tightness at 0.35 and
uses coarse generator proxies for storage and network conditions. Actual
size, saturation indicators and missing-proxy flags should therefore remain
visible alongside an aggregate; none has yet been shown predictive for the
new request-specific outcomes.

\subsection{Compare progressively richer parameterizations}
The proposed hierarchy begins with size and budget, then adds the common
scalar $I$, the separate stress/hardness axes, an interpretable component
subset and all 22 components. Each level is compared with and without
family controls, and against the simple request-conditioned family baseline.
The extra complexity is retained only if it improves grouped validation
loss and the declared routing objective beyond the simpler descriptions.

A candidate joint-success model makes method and deadline interactions
explicit, for example
\begin{equation}
 \operatorname{logit}\widehat p_m
 =\beta_{0m}+f_m(\log d)+\boldsymbol\beta_m^\top\boldsymbol z_s
 +\boldsymbol\gamma_m^\top\boldsymbol z_s\log d
 +\boldsymbol\eta_m^\top h(\tau,\bar E).
 \label{eq:component_router}
\end{equation}
This is a proposed regularized baseline, not a fitted result. The encoding
$h$ uses active-condition indicators and finite threshold values; disabled
constraints are not represented by infinite numeric features. Risk tolerance
and priority act in the selection rule. More flexible component--requirement
interactions must earn their complexity through nested grouped validation.
For a fixed stored run, relaxing a quality or energy cap cannot reduce
acceptance; model coherence should respect these nested events. Independent
budget runs do not justify enforcing monotonicity in budget without testing it.

An optional compact alternative learns a request-conditioned index
$I_{\theta(r)}$ with nonnegative simplex weights, shrinkage toward the common
weights and a small declared number of profiles. This produces a supervised
selection feature, not a newly validated universal physical criticality
measure. Its weights are learned once on development folds and frozen;
there is no scenario-specific retraining or online weight search. Compare
it against keeping $S,H$ or the components directly, since compressing
opposing mechanisms into one scalar can lose useful information.

\subsection{Mechanistic hypotheses are not selection thresholds}
Three groups motivate tests. Size, binary-count and thermal/ramp proxies
may distinguish long-search regimes; storage adequacy, residual variability
and state-of-charge coupling may help identify the VRE cases where JEPA
adds joint coverage; network/trade proxies may distinguish service-preserving
RH responses from economically precise direct solutions. These are
domain-informed hypotheses. The family results do not establish component
effects, signs, useful cutoffs, or a causal explanation for a method's wins.

The same input can affect several components in opposite directions.
For example, greater intertie density decreases the inverse-congestion
stress proxy but increases the interconnection-density hardness proxy.
Thus an action cannot be assigned by a rule such as ``high index implies RH''
without conditioning on deadline, economic tolerance and service cap.
Paper 1's sensitivity diagnostics should identify stable descriptors and
aggregation ambiguities; paper 3 must test their incremental decision value.

\subsection{Nested tuning, stability and deployment safeguards}
The proposed development study varies $\lambda$, simplex weights and
predeclared normalization alternatives inside training folds only. It
compares score/rank changes with paired changes in out-of-fold joint-success
loss, selected-method frequency, request satisfaction and conditional
cost/time regret. A large score change with unchanged decisions has no
demonstrated routing value. A ranking improvement can also be useless if
it does not distinguish the few cases with an alternative acceptable method.

All budgets, repetitions, threshold profiles and variants of one scenario
remain in one group. Reweighting, subset selection, bin boundaries and risk
calibration occur inside the development procedure; the independent test
cannot choose them. No sign constraint on solver success follows solely
from a component's ``stress'' name. Unsupported or highly saturated inputs
are reported as out-of-support warnings and use the declared fallback,
not an invented reliability guarantee. Feature extraction and any extra
computations remain inside the online deadline.

\pending{No sensitivity sweep, feature-importance ranking, fitted weight,
component-conditioned probability or threshold policy is measured in this
revision. The arithmetic analysis in paper 1 and the retrospective ceilings
here specify what to test; they do not provide trained router parameters.}

\section{Candidate models and grouped development}
The proposed first models are deliberately simple: one regularized logistic
or shallow-tree joint-success predictor per method, with cost and exit-time
regressors fitted on the corresponding acceptable outputs. Their inputs
include size descriptors, the 22 common criticality components, execution
allowance, economic tolerance and energy cap. The request priority and risk
tolerance act in the selection rule. Sparse conditional supports must be
reported; an elaborate model is not justified by the existence of four
actions alone.

Baselines include direct HiGHS, the development-selected fixed method for
each request, and request-conditioned family or criticality-bin rules.
Unknown groups use a declared fixed fallback. Family labels provide a useful
descriptive baseline but may not exist in a deployment distribution.
Calibration and model choice are performed within grouped development,
with no privileged reference feature.

All methods, budgets, repetitions and request profiles derived from the same
scenario stay in one group. Paired structural variants share their base
group. A scenario copied across several energy caps or cost thresholds is
not several independent observations. Preprocessing and calibration are
fitted only on training folds. The profiles were developed while
inspecting the earlier subset and the complete 600-case campaign. This entire
cohort remains exploratory; refitting or cross-validating a selector does not
turn it into an untouched confirmatory test.

Out-of-fold joint coverage, probability quality and cost/time contrasts inform
model choice. They are not an unbiased final score for the selected model.
The earlier broader design provides a separate 120-case final cohort and a
controlled structural suite, whose use must preserve group independence.
A design lock must record the exact profiles, caps, zero rule, risk thresholds,
fallback, continuation criteria, selected models and case manifests before
final outcomes are inspected. No new fitting or final evaluation is performed
here.

\section{Online execution and the available budget}
Online feature extraction and selection consume the same decision budget as
the chosen method. Let $T_X$, $T_\pi$ and $T_m$ denote their respective
durations, with method execution including validation, output writing and
closure. The available method allowance is
\begin{equation}
 d=b-T_X-T_\pi .
\end{equation}
Selection must score methods against the remaining allowance, using measured
feature time and a declared conservative estimate of selection overhead
before the choice, then deducting the actual elapsed overhead at dispatch.
A nonpositive remainder is a deadline failure. Scoring with $b$ as if the
selector were free overstates the available computation.

The full decision is evaluated by
\begin{equation}
 Y^\pi_{sr}=\ind\{\text{valid returned output},\
 T_X+T_\pi+T_m\leq b,\
 Q^\pi_s(\tau)=1,\
 E^\pi_s\leq\bar E\},
 \label{eq:online}
\end{equation}
with the same inactive-constraint conventions. The full time also enters the
time-priority comparison. The finalization-reserve policy and the requirement
that the whole process tree terminate remain unchanged.

A lookup of the selected member's full-budget historical outcome cannot
establish Eq.~\eqref{eq:online}. Neither can a 120-second execution be
substituted for an unobserved thirty-second allocation. Retrospective
complementarity, predictive recovery and measured online benefit must remain
separate claims.

The existing architecture permits static library and model preload before
receiving the scenario. Any scenario-specific graph, embedding or feature
computation remains inside the clock. Preloading a portfolio also has memory
and setup costs, which are reported even if excluded from the decision
interval. Unused workers must still close within the deadline. Routing is
not a cost-free table lookup.

\section{Evaluation and conditional reporting}
\subsection{What is established and what remains to be tested}
The final primary records establish the descriptive ceilings in
Table~\ref{tab:final_oracles}, the paired timing contrasts and their common
supports, and the ENS-cap sensitivity. These are observed, retrospective
properties of this portfolio. They neither establish predictive recovery of
the oracle nor guarantee a requested risk level.

\pending{Fit and calibrate candidate selectors only after locking the request
profiles and grouped development procedure. Compare against request-conditioned
fixed and family rules and a HiGHS/RH-only portfolio, include online overhead,
and evaluate on independent cases. Criticality sensitivity, learned weights
and component ablations described here have not been executed.}
The amended service-v2 implementation prepares the five-deadline grid,
request-specific continuation and model choice, grouped calibration,
conditional cost/time prediction and constrained selection. This amendment
has not been executed or validated; its diagnostic output is expected at:
\tablepath{../runs/v2_service_by_deadline/artifacts/routing_diagnostic/tables/oracle_gain.csv}

\subsection{Requirement satisfaction and useful conditional gains}
For each profile--deadline pair, report complete observed support, joint satisfaction,
separate deadline/validity/quality/service violation indicators, fallback
frequency and uncertainty. Violations can overlap; they are not automatically
disjoint failure classes. Cost and total-time differences on common acceptable
outputs are presented with paired support and the all-trial coverage
comparison, so that an attractive successful subset cannot conceal failures.
Calibration and observed violations around the requested risk level must be
shown rather than inferred from mean classification accuracy.

A timed HiGHS-only selection rule separates routing overhead from direct
HiGHS execution. Fixed, family, criticality-bin and learned rules must use the
same request and execution boundary. The primary practical gain is the
profile's declared coverage or cost/time improvement, not an unqualified
five-percent success score or a change in selected-method frequencies.

\pending{Selector comparison, model receipts and timed final results remain
pending. The service-v2 code prepares these outputs; no fitted result or
measured online benefit is claimed here.}
\tablepath{../runs/v2_service_by_deadline/routing/development_predictions.csv}
\tablepath{../runs/v2_service_by_deadline/artifacts/final_combined/tables/joint_paired.csv}
\tablepath{../runs/v2_service_by_deadline/artifacts/final_router/tables/timed_request_rows.csv}

\subsection{Independent and structural scope}
The separate final design contains twenty fresh scenarios per family.
The controlled suite contains thirty bases with paired size, line-structure,
fleet and numbering changes. Claims require uncertainty at the base-scenario
level, with the service profiles preserved across conditions. These are
planned evaluations, not outcomes of the completed historical campaign.

The current action set retains the compatible learned models. Longer
48- and 72-hour cases are outside that set and are not encoded as learned
method failures; a separately trained variable-horizon system belongs to
the companion comparison and would define a different portfolio.
\tablepath{../runs/v2_service_by_deadline/artifacts/controlled_test_combined/tables/joint_coverage.csv}

\section{Discussion}
The relevant choice is conditional on a request rather than an intrinsic
ranking of algorithms. Lower economic cost does not ensure lower unserved
energy, a fast exit does not ensure adequate quality, and an acceptable
output does not certify optimality. The proposed inputs make these distinctions
visible to the decision rule instead of hiding them in one utility weight.

There are several legitimate negative outcomes. A fixed method may already
satisfy a request well enough to leave little useful room. The portfolio may
be complementary but its favorable cases may be hard to recognize before
solving. Sparse or shifted data may prevent reliable probability estimates.
Finally, selection overhead can erase an offline advantage. These outcomes
require different explanations and are not interchangeable with low
classification accuracy.

Risk-conditioned selection is a modeling proposal, not a guaranteed service
policy. A fallback can identify the least unsupported action without making
the requested risk level true. Similarly, a conditional cost or time model
must not conceal failures behind a lower average on a smaller successful
subset. The completed primary grid has no missing portfolio cells
or references; this does not remove synthetic-family selection, imperfect
reference optimality, limited repetitions or the single-hardware scope.
Extensions with missing outcomes must report their support explicitly.

A pipeline member need not win to make the comparison informative. If it
remains dominated for the declared profiles, that is a supported reason not
to select it, rather than a reason to invent a specialty. The proposed
study must allow RH, HiGHS or another fixed choice to remain preferable
where the evidence warrants it.

\section{Conclusion}
The completed comparison demonstrates request-dependent complementarity and
quantifies its ceilings. Short-budget VRE-with-batteries cases provide a
concrete joint-coverage opportunity, while a simple family-conditioned rule
already captures most continuity coverage at thirty seconds. Cost-only
five-percent coverage at 120 seconds leaves no gain over HiGHS; useful
conditional time improvements remain possible but are modest in aggregate
before paying selection overhead.

Criticality supplies candidate explanatory inputs, not an already validated
action map. Its score sensitivity, predictive value beyond size and family,
and effect on selected decisions must be tested separately. A demand-dependent
router would estimate joint success, optimize the stated priority among
supported methods and flag unsupported requests. Whether this recovers a
useful fraction of the retrospective opportunity with calibrated reliability
and positive online benefit remains open.

\section*{Provenance and execution status}
The final campaign completed on 10 November 2026. Its accepted analysis
contains 16,800 primary executions on 600 scenarios and 1,680 additional
executions on 30 cases. All primary method--budget cells and all 600 references
are available; 383 references are certified optimal. The final acceptance
file reports no invalid or late output counted as a success. Failed trials
remain in the denominators.

The service-aware descriptive analysis was produced on 11 November 2026
before this manuscript revision. An unchanged copy is archived at
\path{notes/final_results_2026-11-11.json}; it contains source-file hashes,
family and budget summaries, paired supports, retrospective ceilings,
repetition summaries, ENS-cap sensitivity and conditional timing contrasts.
Its source directory, relative to the repository root, is
\path{outputs/thermal_transitions_v1/operational_exit_v1/analysis_service_profiles/20261111_143949/}.
The accepted campaign tables are under
\path{outputs/thermal_transitions_v1/operational_exit_v1/analysis/}.
The earlier \path{notes/resultats_fournis_2026-11-06.txt} is retained as
historical context and is superseded for numerical claims in these drafts.

This writing revision reuses those existing results and inspects index source
definitions. It launches no new experiment, solver execution, scenario
generation, training, parameter sweep or empirical index analysis. It does not
change runtime code, model weights, frozen campaign records or bibliographies.
Textual consistency checks are separate from \LaTeX{} compilation and rendered
PDF inspection, neither of which is performed in this revision.

No requirement-aware router, learned criticality weighting, sensitivity grid
or component-ablation study is trained or run in this writing revision.
The service-v2 code now prepares the request-aware analyses and router.
Learned criticality weights and the full sensitivity/subset study remain
separate proposals; no training or new experiment was run in this update.

\bibliographystyle{unsrtnat}\bibliography{paper3/references}
\end{document}
